\section{Devices and models}
\label{sec:models}

\subsection{Geometries}

All devices are two-terminal planar channels with the source contact on the
left and the drain on the right. The funnel [Fig.~\ref{fig:devices}(a)] has
the dimensions of device D of Nguyen \textit{et al.}~\cite{nguyen2021geometric}:
leads $2\um$ wide, a taper at 30 degrees to the axis down to a neck of
$175\nm$, and an abrupt shoulder on the drain side; its measured asymmetry is
$1.66$ at $80\ \mu$A. The classical Tesla valve [Fig.~\ref{fig:devices}(b)] is
built from channel centre lines: a zig-zag main channel of width $175\nm$
inclined at $\pm 25$ degrees, and at every vertex a loop that continues the
arriving segment straight on, turns through a half circle of radius $175\nm$
and rejoins the main channel against the stream. Flow from left to right
follows the zig-zag; flow from right to left is split into the four loops. We
also keep three earlier variants in which the loops are attached to one side
of a straight channel, because they illustrate that the conclusion does not
depend on the details of the valve, and a straight $175\nm$ channel between
the same leads as a control. For the terahertz calculations each shape is
also scaled so that its narrowest width is $1\um$, with $4\um$ contacts,
which brings the device resistance from kilohms to the hundred-ohm range of
an antenna [Fig.~\ref{fig:devices}(c)].

\subsection{Ballistic model}
\label{sec:model-ballistic}

The ballistic model is the semiclassical trajectory (Monte
Carlo~\cite{jacoboni1983monte}) method of the geometric-diode
literature~\cite{zhu2013graphene,stearns2021simulation,truccolo2024comprehensive}.
Electrons with the Dirac dispersion $E=\hbar\vF|\bm{k}|$,
$\vF=10^{6}$~m/s, are injected at the contacts with the flux-weighted cosine
law, fly along classical trajectories, reflect from the walls specularly with
probability $p$ and diffusely otherwise, and suffer isotropic
momentum-relaxing collisions after exponentially distributed free paths of
mean $\lambda$. The current is the Landauer expression
\begin{equation}
I=\frac{g e}{h}\int dE\,N(E)\,T(E;V)\,\bigl[f_{S}(E)-f_{D}(E)\bigr],
\label{eq:landauer}
\end{equation}
with $g=4$, $N(E)=EW/\pi\hbar\vF$ the number of modes of the injecting contact
of width $W$, $T$ the transmission probability estimated from the
trajectories, and $f_{S,D}$ the Fermi functions of the contacts at
electrochemical potentials $\EF$ and $\EF-eV$. Energies are sampled from the
bias window $N|f_{S}-f_{D}|$, so the estimator has zero variance at zero bias.

With no potential inside the device, flux conservation gives
$N_{S}T_{SD}=N_{D}T_{DS}$ at every energy and Eq.~\eqref{eq:landauer} is odd
in $V$ for any shape; the trajectories cannot rectify. Rectification enters
through the potential energy $U(\bm{r};V)$ that the bias creates inside the
channel~\cite{christen1996gauge}. We use the simplest assumption of the published
simulations, a uniform field along the channel, $U=-eVx/L$ over the device
length $L$, and integrate the curved flights in chords with a controlled
sagitta. Self-consistent electrostatics is known to roughly halve the
asymmetry of such diodes relative to the uniform-field value and to
reverse it for shallow tapers~\cite{truccolo2024comprehensive}; we return
to the field distribution in Sec.~\ref{sec:ballistic}. The figure of merit is the asymmetry $A=I(+V)/|I(-V)|$; writing
$I=G_{1}V+G_{2}V^{2}$, $A\simeq 1+2(G_{2}/G_{1})V$, and a cosine drive of
amplitude $V_{a}$ rectifies $\Idc=G_{2}V_{a}^{2}/2$ if the response is
instantaneous. Carriers are electrons at density $n=1.4\times10^{12}\percm$
($\EF=0.138$~eV) and $300$~K, the density we estimate from the gate voltage
and lever arm reported for that device~\cite{nguyen2021geometric}; holes
give the inverse asymmetry.

The solver was validated on eleven tests before use: unit transmission of a
specular channel, the free-molecular transmission factors of a diffuse
channel~\cite{richley1964numerical}, the Sharvin
conductance~\cite{sharvin1965possible}, the ballistic-to-diffusive crossover, reciprocity with and
without the field, unit asymmetry of a mirror-symmetric constriction, and
convergence of the chord integration. Classical trajectories require
features well above the Fermi wavelength, $30\nm$ here.

\subsection{Viscous model}
\label{sec:model-viscous}

Where electron-electron collisions are the fastest scattering process the
electron liquid obeys the Navier-Stokes
equation~\cite{gurzhi1968hydrodynamic,bandurin2016negative}. The dc model is
steady incompressible flow with momentum relaxation,
\begin{equation}
\Rey\,(\bm{u}\cdot\nabla)\bm{u}=-\nabla p+\nabla^{2}\bm{u}-\Da\,\bm{u},
\qquad \nabla\cdot\bm{u}=0,
\label{eq:ns}
\end{equation}
in units of the lead width $W$ and the mean inlet speed, with
\begin{equation}
\Rey=\frac{I}{n e \nu},\qquad \Da=\frac{W^{2}}{\nu\taumr},\qquad
S=\frac{\ell_{s}}{W}
\end{equation}
the Reynolds number, the damping (Gurzhi) number and the slip length
$\ell_{s}$ of the Navier wall condition~\cite{kiselev2019boundary} in units
of $W$. In two dimensions the Reynolds number depends on the current and not
on the width. The device is current biased with a developed inlet profile,
and the figure of merit is the diodicity $D=\Pi_{r}/\Pi_{f}$, the ratio of
the pressure (voltage) drops for flow from drain to source and from source
to drain. Equations~\eqref{eq:ns} are solved with Taylor-Hood
elements~\cite{taylor1973numerical} on triangular
meshes~\cite{geuzaine2009gmsh} with Newton iteration and a sparse direct
solver~\cite{guennebaud2010eigen}; the wall condition
$\bm{u}\cdot\bm{n}=0$ is imposed with length-weighted
normals~\cite{engelman1982implementation}, which makes the discrete wall flux
vanish exactly and removes a spurious diodicity of $10^{-5}$ that polygonal
walls otherwise produce. Validation covers Poiseuille and Brinkman flow with
slip, Stokes reciprocity ($D-1<10^{-10}$ at $\Rey=0$), the Kovasznay
solution~\cite{kovasznay1948laminar}, the backward-facing
step~\cite{armaly1983experimental,erturk2008numerical}, and the
diodicity of the valve itself as a classical fluid device. The discretisation
floor of $D-1$ is $3\times10^{-5}$.

\subsection{Plasmonic model}
\label{sec:model-plasmonic}

At terahertz frequencies the fluid is compressible: a gate at distance $d$
makes the local density follow the local electrochemical potential, and
density waves propagate as plasmons. With the gradual-channel approximation
and degenerate Dirac carriers the electrochemical potential is
$\Phi(n)=e^{2}n/C_{g}+\hbar\vF\sqrt{\pi n}$, the hydrodynamic mass is
$m^{*}(n)=\hbar\sqrt{\pi n}/\vF$~\cite{tomadin2013theory,svintsov2012hydrodynamic}, and because both
depend on the density only the force per unit mass is the gradient of a
single function $h(n)$ with $h'=\Phi'/m^{*}$. The model is
\begin{align}
\partial_{t}n+\nabla\cdot(n\bm{u})&=0,\label{eq:cont}\\
\partial_{t}\bm{u}+(\bm{u}\cdot\nabla)\bm{u}+\frac{\bm{u}}{\taumr}
&=-\nabla h(n)+\nu\nabla^{2}\bm{u},\label{eq:mom}
\end{align}
with the plasma speed $s(n)$ given by
\begin{equation}
s^{2}(n)=n\,h'(n)=\frac{e^{2}n}{C_{g}m^{*}(n)}+\frac{\vF^{2}}{2}.
\label{eq:s2}
\end{equation}
The second term is the quantum capacitance~\cite{fang2007carrier}; its share $q=\vF^{2}/2s_{0}^{2}$
of $s_{0}^{2}=s^{2}(n_{0})$ is $0.04$ at $n_{0}=10^{12}\percm$ under $30\nm$ of
hBN and grows to $0.5$ at the thin-gate, low-density corner of
Sec.~\ref{sec:operating}. Replacing $h$ by $s_{0}^{2}n/n_{0}$ gives the
shallow-water model of Dyakonov and Shur~\cite{dyakonov1993shallow}. At
frequencies above $1/\tauee$ the viscosity is taken in the viscoelastic form
$\nu/(1+i\omega\tauee)$.

The antenna is a Thevenin source of amplitude $V_{a}$ and resistance $\Rant$
connected to lumped ports: each contact is an equipotential whose voltage and
current satisfy Kirchhoff's law with the external circuit. Two terminations
are used throughout. In the \emph{symmetric} one the antenna drives the two
contacts in antiphase through $\Rant/2$ each and the gate sits at the
midpoint; the circuit is then mirror symmetric and any dc response is due to
the shape of the channel. In the \emph{Dyakonov-Shur} (DS) termination the
antenna drives the source through $\Rant$ and the drain is tied to the gate;
the terminations are then asymmetric and a uniform channel already
rectifies~\cite{dyakonov1996detection}. The responsivity is the
short-circuit dc current per available power,
\begin{equation}
\RI=\frac{I_{\mathrm{dc}}}{\Pav},\qquad \Pav=\frac{V_{a}^{2}}{8\Rant},
\end{equation}
and the Johnson-noise-limited noise-equivalent power is
$\mathrm{NEP}=\sqrt{4\kB T/\Rdc}/\RI$ with $\Rdc$ the dc resistance of the
device.

Two solvers are used. For small signals the fields are expanded to second
order about rest: the first order is a complex time-harmonic plasmon problem,
and the time average of the second order is a Stokes-Brinkman problem for the
dc flow, forced by products of first-order fields, whose port current or
voltage is the rectified signal (Appendix~\ref{app:perturbation}). For finite amplitudes
Eqs.~\eqref{eq:cont} and \eqref{eq:mom} are integrated in time with the
second-order backward difference formula and Newton iteration, with a bias
tee that separates the dc load from the antenna. Both use Taylor-Hood
elements on the meshes of the viscous model. The validation ladder has
seventeen tests: the static limit against the viscous solver, the plasmon
dispersion, the Dyakonov-Shur closed form [Fig.~\ref{fig:small}(a)],
reciprocity of the symmetric termination, mesh and time-step convergence,
agreement of the two solvers at small amplitude, and invariance under
co-scaling of the physical parameters at fixed nondimensional groups.

\begin{table}
\caption{Parameter sets. The ballistic model follows the experiment of
Ref.~\cite{nguyen2021geometric}; the two plasmonic sets are representative of
hBN-encapsulated graphene at the two temperatures and are labelled by them;
they are assumptions, not fits. The $100$~K viscosity is of the order
measured~\cite{bandurin2016negative,krishnakumar2017superballistic}; the
$300$~K value is the kinetic estimate $\vF^{2}\tauee/4$. $W$ is the contact
width.}
\label{tab:params}
\begin{ruledtabular}
\begin{tabular}{ll}
Quantity & Value \\
\hline
\multicolumn{2}{l}{\textit{Ballistic}} \\
$n$, $\EF$, $T$ & $1.4\times10^{12}\percm$, $0.138$ eV, 300 K \\
specularity $p$ & 1, 0.5, 0 \\
mean free path $\lambda$ & $\infty$, 700, 200 nm \\
\multicolumn{2}{l}{\textit{Viscous}} \\
current at $\Rey=1$ & $160\ \mu$A ($\nu=0.1$ m$^2$/s), $32\ \mu$A (0.02) \\
$\Da$, $S$ & 0 to 128; 0 to $\infty$ \\
\multicolumn{2}{l}{\textit{Plasmonic}} \\
$n_{0}$, $d$, $\epsilon_{r}$ & $10^{12}\percm$, 30 nm, 3.8 \\
$s_{0}$, gate overdrive & $3.6\,\vF$, 1.43 V \\
``100 K'': $\taumr$, $\tauee$, $\nu$ & 5 ps, 1 ps, 0.1 m$^2$/s \\
``300 K'': $\taumr$, $\tauee$, $\nu$ & 1 ps, 0.1 ps, 0.03 m$^2$/s \\
$W/s_{0}\taumr$ ($W=4\um$) & 0.22 (100 K), 1.12 (300 K) \\
antenna $\Rant$ & $75\ \Omega$ \\
\end{tabular}
\end{ruledtabular}
\end{table}
